% ============================================================
%  Příprava na maturitu – Mechatronikaa
% ============================================================

\documentclass[a4paper,10pt]{article}

\usepackage[T1]{fontenc}
\usepackage[utf8]{inputenc}
\usepackage[czech]{babel}
\usepackage[left=2.4cm, right=1.8cm, top=1.8cm, bottom=1.8cm]{geometry}
\usepackage{lmodern}
\usepackage{microtype}
\usepackage{parskip}
\usepackage{titlesec}
\usepackage{xcolor}
\usepackage{hyperref}
\usepackage{tabularx}
\usepackage{booktabs}
\usepackage{tikz}
\usepackage{mdframed}
\usepackage{amsmath}

\usetikzlibrary{arrows.meta, positioning}

% --- Barvy ---
\definecolor{accent}{HTML}{03254E}
\definecolor{lightblue}{HTML}{EAF2F8}
\definecolor{lightyellow}{HTML}{FDFBE4}
\definecolor{lightgreen}{HTML}{EAF7EA}

% --- Boxy ---
\newmdenv[backgroundcolor=lightblue,linecolor=accent!40,linewidth=1pt,roundcorner=4pt]{pojmybox}
\newmdenv[backgroundcolor=lightyellow,linecolor=accent!30,linewidth=1pt,roundcorner=4pt]{vysvetlenibox}
\newmdenv[backgroundcolor=lightgreen,linecolor=accent!30,linewidth=1pt,roundcorner=4pt]{vzorcebox}

\begin{document}

% ============================================================
\section{Základy PLC a mikrokontrolérů}

\begin{pojmybox}
PLC, scan cyklus, I/O, CPU, ladder diagram,
mikrokontrolér, GPIO, UART, SPI, I2C,
interrupt, timer, PWM, register, firmware
\end{pojmybox}

% ============================================================
\begin{vysvetlenibox}

\textbf{PLC (Programmable Logic Controller)}\\
PLC je průmyslový řídicí systém navržený pro spolehlivý a dlouhodobý provoz.

\textbf{Hlavní vlastnosti:}
\begin{itemize}
\item deterministické řízení (stejný čas cyklu)
\item odolnost vůči rušení (EMC)
\item modulární struktura
\item jednoduchá údržba
\end{itemize}

% -----------------------------
\textbf{Scan cyklus PLC}

\begin{center}
\begin{tikzpicture}[
node/.style={rectangle, draw=accent, fill=lightblue,
rounded corners=3pt, minimum width=3.2cm, minimum height=0.6cm, font=\small}
]

\node[node] (a) {1. Čtení vstupů};
\node[node, below=of a] (b) {2. Vykonání programu};
\node[node, below=of b] (c) {3. Zápis výstupů};
\node[node, below=of c] (d) {4. Diagnostika};

\draw[->] (a)--(b);
\draw[->] (b)--(c);
\draw[->] (c)--(d);
\draw[->] (d.west) -- ++(-1,0) -- (a.west);

\node[right=0.3cm of b, font=\scriptsize] {typicky 1–10 ms};

\end{tikzpicture}
\end{center}

 PLC nepracuje „real-time per event“, ale cyklicky.

\textbf{Programovací jazyky PLC (IEC 61131-3)}

PLC se programují podle normy IEC 61131-3, která definuje několik jazyků:

\begin{itemize}
\item \textbf{LD (Ladder Diagram)} – žebříčkový diagram
\item \textbf{FBD (Function Block Diagram)} – blokové schéma
\item \textbf{ST (Structured Text)} – textový jazyk podobný Pascalu
\item \textbf{IL (Instruction List)} – instrukční seznam (historický)
\item \textbf{SFC (Sequential Function Chart)} – sekvenční grafy
\end{itemize}

% -----------------------------
\textbf{LD – Ladder Diagram}

Nejpoužívanější v průmyslu. Vypadá jako elektrické schéma relé.

\begin{itemize}
\item levá strana = vstupy
\item pravá strana = výstupy
\item logika AND/OR pomocí kontaktů
\end{itemize}

Výhoda: snadné pochopení pro elektrikáře.

% -----------------------------
\textbf{ST – Structured Text}

Textový jazyk podobný C/Pascalu:

\begin{vzorcebox}
\begin{verbatim}
IF temperature > 50 THEN
    motor := TRUE;
ELSE
    motor := FALSE;
END_IF;
\end{verbatim}
\end{vzorcebox}

Výhoda:
\begin{itemize}
\item složitější algoritmy
\item matematika
\item PID regulace
\end{itemize}

% -----------------------------
\textbf{FBD – Function Block Diagram}

Grafické bloky:

\begin{itemize}
\item AND / OR / NOT bloky
\item funkční bloky (timer, counter)
\item propojení signálů mezi bloky
\end{itemize}

Výhoda: přehlednost u řízení procesů.

% -----------------------------
\textbf{SFC – Sequential Function Chart}

Řídí sekvence kroků:

\begin{itemize}
\item kroky (steps)
\item přechody (transitions)
\item paralelní větvení
\end{itemize}

Použití:
\begin{itemize}
\item výrobní linky
\item automatické cykly
\end{itemize}

% -----------------------------
\textbf{IL – Instruction List (zastaralý)}

Nízká úroveň, podobná assembleru.

Dnes už se prakticky nepoužívá (deprecated).


% ============================================================
\textbf{Mikrokontrolér (MCU)}\\
Malý počítač na jednom čipu.

\textbf{Obsahuje:}
\begin{itemize}
\item CPU
\item RAM / Flash
\item GPIO piny
\item periferie (UART, SPI, I2C, ADC, timers)
\end{itemize}

% ============================================================
\textbf{PLC vs mikrokontrolér}

\begin{tabularx}{\linewidth}{l X X}
\toprule
 & PLC & Mikrokontrolér \\
\midrule
Cíl & průmysl & vestavěné systémy \\
Cena & vysoká & nízká \\
Odolnost & velmi vysoká & střední \\
Programování & IEC 61131-3 & C / C++ \\
Flexibilita & omezená & vysoká \\
Reakce & cyklická & event-driven \\
\bottomrule
\end{tabularx}

% ============================================================
\textbf{GPIO (mikrokontrolér)}

\begin{itemize}
\item vstup = čtení senzoru
\item výstup = řízení LED, motoru
\item alternativní funkce = UART/SPI/I2C/PWM
\end{itemize}

% ============================================================
\textbf{Interrupts (přerušení)}

Přerušení umožňuje okamžitou reakci na událost.

\begin{itemize}
\item tlačítko
\item časovač
\item UART příjem dat
\item změna signálu
\end{itemize}

 Výhoda: nemusí se čekat v loop()

% ============================================================
\textbf{Timers (časovače)}

Použití:
\begin{itemize}
\item měření času
\item PWM generace
\item periodické úlohy
\end{itemize}

\[
f_{out} = \frac{f_{cpu}}{prescaler}
\]

% ============================================================
\textbf{PWM}

\begin{vzorcebox}
\[
D = \frac{t_{on}}{T} \cdot 100\,\%
\]
\end{vzorcebox}

Použití:
\begin{itemize}
\item řízení motorů
\item LED stmívání
\item regulace výkonu
\end{itemize}

% ============================================================
\textbf{Komunikace MCU}

\begin{itemize}
\item UART – jednoduchá sériová linka
\item SPI – rychlá synchronní komunikace
\item I2C – více zařízení na 2 vodičích
\end{itemize}

% ============================================================
\textbf{Firmware}

Firmware = program v mikrokontroléru:
\begin{itemize}
\item běží přímo na HW
\item často bez OS
\item musí být stabilní
\end{itemize}

% ============================================================
\end{vysvetlenibox}

\end{document}
